\documentclass[aspectratio=169]{beamer}
\usetheme{Madrid}
\usecolortheme{default}

% Packages
\usepackage{graphicx}
\usepackage{booktabs}
\usepackage{amsmath}
\usepackage{hyperref}
\usepackage{xcolor}

\title{CARE V4: Analyzing the ResNet Bottleneck}
\subtitle{Discovering the BatchNorm Nullification Effect}
\author{AUB Research Team}
\date{\today}

\begin{document}

\begin{frame}
    \titlepage
\end{frame}

\begin{frame}{Recap: The Original Promise of CARE}
    \textbf{The Hypothesis:}
    \begin{itemize}
        \item standard surrogate gradients lack mechanisms to recover dead/silent pathways in deep Spiking Neural Networks (SNNs).
        \item CARE (Content-Aware Re-Activation) applies homeostatic magnitude scaling to incoming weights of silent neurons, acting as a dynamic "capacity insurance."
    \end{itemize}
    
    \vspace{0.3cm}
    \textbf{Early Success in VGGs:}
    \begin{itemize}
        \item In VGG architectures lacking residual connections, CARE successfully rescued models experiencing catastrophic starvation, reducing dead neurons from ~99\% to ~1\%.
    \end{itemize}
\end{frame}

\begin{frame}{The Scaling Failure: Translating to ResNet}
    \textbf{The NeurIPS Rigorous Ablation Suite Results:}
    \begin{itemize}
        \item Moving to ResNet-18/34 block structures exposed a fundamental issue.
        \item We observed a significant performance "Tax" instead of enhanced capacity.
    \end{itemize}
    
    \vspace{0.3cm}
    \begin{table}[]
        \centering
        \begin{tabular}{lrrcc}
            \toprule
            Setting & Dataset & Control & CARE & $\Delta$ Accuracy \\
            \midrule
            Phase 1 & CIFAR-100 & \textbf{38.5\%} & 23.8\% & -14.7\% \\
            Phase 2 & CIFAR-10 ($T=2$) & \textbf{66.5\%} & 45.2\% & -21.3\% \\
            \bottomrule
        \end{tabular}
        \caption{NeurIPS ResNet-18 Baseline Suite (30 Epochs)}
    \end{table}

    \begin{itemize}
        \item \textbf{Conclusion:} CARE introduces detrimental weight updates without the intended recovery benefits in properly initialized ResNets.
    \end{itemize}
\end{frame}

\begin{frame}{The Starvation Test Reality}
    \textbf{Testing the "Absolute Rescue" Stress Case:}
    \begin{itemize}
        \item We re-ran the Starvation Stress Test using high weight decay ($\sigma = 5 \times 10^{-3}$) to force capacity collapse in ResNet-18 (CIFAR-10).
    \end{itemize}
    
    \vspace{0.2cm}
    \begin{columns}
        \begin{column}{0.5\textwidth}
            \begin{block}{Static Solution (Control)}
                \begin{itemize}
                    \item \textbf{Accuracy:} 25.4\%
                    \item \textbf{Dead Ratio:} 0\%
                    \item \textbf{Status:} Survived easily due to normal initialization and residual paths.
                \end{itemize}
            \end{block}
        \end{column}
        \begin{column}{0.5\textwidth}
            \begin{block}{Dynamic Solution (CARE)}
                \begin{itemize}
                    \item \textbf{Accuracy:} 22.8\%
                    \item \textbf{Dead Ratio:} 0.3\%
                    \item \textbf{Status:} Underperformed the control group.
                \end{itemize}
            \end{block}
        \end{column}
    \end{columns}

    \vspace{0.4cm}
    \centering
    \textit{Finding: The Control ResNet does not fail catastrophically under extreme decay. The "Absolute Rescue" hypothesis does not hold for properly initialized Residual Networks.}
\end{frame}

\begin{frame}{The Breakthrough: The BatchNorm Nullspace}
    \textbf{Why Does Homeostasis Fail in ResNets?}
    \begin{itemize}
        \item Almost all modern deep SNNs (including our SEW-ResNets) use Batch Normalization (BN) prior to the activation functions.
        \item CARE's primary mechanism involves scaling weight \textbf{magnitudes} to increase pre-synaptic variance until thresholding occurs.
    \end{itemize}
    
    \vspace{0.3cm}
    \begin{block}{The Nullification Effect}
        \begin{enumerate}
            \item CARE increases weight magnitudes $W \leftarrow W + \Delta_w$.
            \item This increases the variance of the convolution output $x_{pre} = W * x$.
            \item BatchNorm mathematically squashes $x_{pre}$ back to a unit variance $\sigma^2 = 1.0$.
            \item The drive to the neuron remains unchanged; \textbf{the intended magnitude boost is perfectly nullified}.
        \end{enumerate}
    \end{block}
\end{frame}

\begin{frame}{The Actual Effect of CARE (The Poisoning Mechanism)}
    \textbf{What exactly is the "Plasticity Tax"?}
    \begin{itemize}
        \item Because the overall variance boost is cancelled out, CARE does not rescue neurons.
        \item However, because weights update asynchronously and asymmetrically depending on their signs, CARE still rotates the effective gradient \textbf{direction}.
        \item This introduces unsupervised noise rotating the learned features away from the task objective without providing the homeostatic bailout.
    \end{itemize}
    
    \vspace{0.3cm}
    \centering
    \textit{CARE acts as pure "Directional Noise" in the presence of Batch Normalization.}
\end{frame}

\begin{frame}{Future Directions and Pivots}
    \textbf{Where do we go from here?}
    
    \begin{itemize}
        \item \textbf{Option A: Targeting the Gamma ($\gamma$):}
        Since BN nullifies weight magnitude but uses its own affine scaled parameter ($\gamma$) to control post-BN variance, we can apply CARE's SNR-gated homeostasis directly to the BN $\gamma$ scale parameters.
        
        \item \textbf{Option B: Alternative Normalization:}
        Investigate whether CARE can safely replace BatchNorm completely (e.g., combining CARE with Weight Standardization), offering a theoretical advantage in sequence-processing SNNs.
        
        \item \textbf{Option C: Shifting the Narrative:}
        Publish a focused analysis paper detailing why standard STDP-like homeostasis fails in SNNs that employ Deep Learning tricks (like BN), highlighting the mathematical nullspace boundary.
    \end{itemize}
\end{frame}

\end{document}
